\subsection*{A stronger local axial-stress margin on the two A.11 supports}
\label{subsec:A11-strong-20261021}
Write $a=1-\lambda$, $b=1+2\lambda$, $T_f=512$ and
$L_p=30\log(1/\lambda)$. Use the actual normalized source
$\sigma'$ bumps $\beta_i$ centered at $L_p-3$ and $L_p-1$.
After the parameter-flattening interpolation the unedited
azimuthal field $E_0$ is independent of $\eta$.
Set $E=E_0(1+\gamma)$,
$\gamma=\sum_{i=1}^{2}c_i(\eta)\beta_i$.
The incoming moment ratio $r_p=I_p/(X_pH_p)$ is independent
of $\eta$ and obeys $0\le r_p\le1/a$.
The exact A.10 entrance formula (proved above) now gives
\[
 0\le r_f\le(a-1/10)^{-1},\qquad
 \sup_{|\eta|\le1}|r_{f,\eta}|<1/700000.
\]
The angular-moment deficit, normalized at the first bump center,
is $\delta=(a^{-1}-r_f)e^{-a(L_p-3)}$. Hence
\begin{equation}\label{eq:A11-strong-delta}
 |\delta|\le50\lambda^{28},\qquad
 |\delta_\eta|\le3\cdot10^{-5}\lambda^{28},
 \qquad0<\lambda\le1/20.
\end{equation}
To retain the actual nonlinear pressure-neutral constraint, define
\[
 u=\int e^{at}\beta(t)\,dt,\quad
 p=\int e^{-bt}\beta(t)\,dt,\quad
 w=\int e^{-bt}\beta(t)^2\,dt,\quad
 x=e^{2a},\quad z=e^{-2b}.
\]
Both exact A.11 equations are
\[
 u(c_1+xc_2)=\delta,\qquad
 2p(c_1+zc_2)+w(c_1^2+zc_2^2)=0.
\]
The source bump satisfies $1/2<u,p<2$, $6<x<8$,
$0<z<1/4$, $w<427$,
$\|\beta\|_\infty\le640/3$ and
$\|\beta'\|_\infty<250000$. Solving the first linear
equation and then the second by contraction on
$\|c\|_\infty\le|\delta|$ gives a map with
\[
 \|\mathcal F(c)\|_\infty
 \le(8/23)|\delta|+743\|c\|_\infty^2,\qquad
 \operatorname{Lip}\mathcal F\le1485|\delta|<1/100.
\]
Thus this ball contains a unique solution, and
differentiation of the \emph{exact} two-dimensional system yields
\[
 \|c\|_\infty\le|\delta|,\quad
 \|c_\eta\|_\infty\le|\delta_\eta|,\quad
 |\gamma|<1/1000,\quad |\gamma_y|<\lambda/4,\quad
 |\partial_\eta\log E|<1/100.
\]
In particular the corrected slope satisfies
$\ell=-\lambda+\partial_y\log(1+\gamma)\le-\lambda/2$.

\begin{theorem}[Improved conditional A.11 axial-stress bound]
\label{thm:A11-support-strong-20261021}
Assume the full source terminal gauges
$S(\infty,\eta)=\Pi(\infty,\eta)=0$ and the prescribed
post-correction stage slopes, including the one-unit
release transition, a $4\log(1/h)$ slope-$-1$ hold,
and $\ell\le-3h/4$ thereafter. Take
\[
 T_f=512,\quad0<\lambda\le
 \min\{1/20,(1+\log P_*)^{-1}\},\qquad
 0<h\le\min\{1/10000,\lambda/100,e^{-T_d}/100\}.
\]
For both full source-corrected A.11 supports,
\begin{equation}\label{eq:A11-strong-support-cone-20261021}
\sup_{\substack{y\in\operatorname{supp}\beta_1
 \cup\operatorname{supp}\beta_2\\|\eta|\le1}}
  \frac{|N_s(y,\eta)|}{E(y,\eta)}
 <\frac{2}{25}\lambda<\frac\lambda2.
\end{equation}
\end{theorem}
\begin{proof}
Because $U=M=J=0$ on these supports, (4.16) and the
genuine terminal gauges give the exact future-energy formula
\[
 N_s=\int_y^\infty\left[
 2\eta(he^{v-y}-A)E(v)^2+
 (1-\eta^2)(1-e^{v-y})E(v)E_\eta(v)\right]dv,
 \qquad A=\tfrac12+h.
\]
Put $\mathfrak d=(2XE^2)^{-1}\int_X^\infty E^2dx$.
The remaining support-to-plateau distance is at most $63/20$.
The release transition, steep hold and final decay imply
$\mathfrak d\le(63/20+1)/2+1/4+h^7/3<3$.
After interpolation the only $\eta$ dependence occurs on the
two bumps, whose combined span is $23/10<3$.
Every later slope is nonpositive, so
\[
 |N_s|/E\le E[4h\mathfrak d+2A+
 3\sup|\partial_\eta\log E|]<2E.
\]
The source reference, pulse and intermediate-stage decay
give, on both corrected supports,
\[
 E<20P_*e^{-13/(2\lambda)}\lambda^{15}.
\]
Since $\lambda\le(1+\log P_*)^{-1}$,
$P_*e^{-13/(2\lambda)}\le e^{-13/2}<1/500$;
therefore $E<\lambda^{15}/25\le\lambda/25$.
Combining the last two inequalities proves the theorem.
\end{proof}

This improves the previous \emph{conditional} A.11 bound
$|N_s|/E<63\lambda/250$ with an explicit sufficient
smallness condition on $\lambda$ relative to $P_*$.
The full global correction parameter order and compatibility
of the terminal gauges remain to be certified; the theorem
does not establish the Appendix-B contraction, unforced
startup solutions or the proposed unforced Navier--Stokes
singularity. The accompanying Lean file proves only finite
scalar inequalities; no complete Lean proof is claimed.
